Phase I is a discovery study: its role is to establish the phenomenon, the
taxonomy, and the existence of task-side headroom, and to motivate the
confirmatory design. None of its results is confirmatory: 24 of its 151
evaluation items overlapped protocol-development measurements, the human
control was designed after inspecting development traces, and its judge was
a conservative proxy (Appendix~\ref{app:phase1-provenance}). We report it
with those limits stated.

\subsection{Harnesses, outcomes, and the three metrics}\label{sec:metrics}
A \textbf{harness} $H$ wraps a frozen target model $M$: it receives the
natural-language task $x$, issues one or more calls to $M$, optionally
executes intermediate artifacts (e.g.\ SQL), and returns a final answer
scored programmatically. With $H_0$ the bare single-call baseline and
$Y \in \{0,1\}$ correctness, the \textbf{harness marginal value} is
$\Delta(x,H,M) = Y(x,H,M) - Y(x,H_0,M) \in \{-1,0,+1\}$
($+1$ repair, $-1$ harness-induced harm).

We propose that harness-generation and
harness-routing pipelines report three precondition metrics on a held-out
set, all computable from the outcome matrix any evaluation already produces:
\textbf{(i) pairwise outcome disagreement} --- over all harness pairs, the
fraction of tasks where correctness differs; \textbf{(ii) union repair
rate} --- the fraction of bare errors some population member repairs;
\textbf{(iii) bare-inclusive oracle headroom}
\[
H(P) \;=\; \mathrm{Acc}\big(x \mapsto \max_{H \in P \cup \{H_0\}} Y(x,H)\big)
\;-\; \max_{H \in P \cup \{H_0\}} \mathrm{Acc}(H).
\]
Positive headroom means population value is \emph{conditional on the item} ---
the signal a router needs. The metrics are diagnostic in combination: high
disagreement with no repair is noise; repair with no headroom means one
harness dominates.

\subsection{The collapse, quantified}\label{sec:collapse}
The domain is BIRD text-to-SQL (SQLite; execution-accuracy judging), target
frozen at \textbf{GLM-5.3-Flash} (temperature 0) throughout Phase I. On the
day-1 population --- 14 harnesses (\bare, \react, and 12 evolution
candidates) $\times$ 60 tasks --- average pairwise outcome disagreement is
\textbf{3.3\%} and \textbf{36 of 91 pairs are identical on every task};
among the 12 generated candidates alone (66 candidate pairs; mean pairwise
code similarity 0.143 under the normalized-character metric of the original
code-difference check), disagreement is 3.8\% with 21 identical pairs.
Every candidate had passed a code-difference check (``all files distinct,
non-trivial diffs''). Union repair over bare errors is 7.7\% and oracle
headroom 3.3\,pp; \react{} differs from \bare{} on zero tasks.

Figure~\ref{fig:phenomenon} makes the central relation quantitative: across
all 91 pairs, pairwise source-code similarity under the figure's token-sequence
metric (mean 0.190; distinct from the candidates-only normalized-character
metric of 0.143 above) has Spearman $\rho = -0.010$ with outcome
disagreement --- code similarity
shows \emph{no observed rank association} with behavioral similarity in this
population. (Pairs share harnesses and are not independent samples; we
report $\rho$ as a descriptive summary, not an inferential test.)

\begin{figure}[t]
\centering
\includegraphics[width=0.78\linewidth]{figures/fig_phenomenon.pdf}
\caption{\textbf{Syntactic diversity vs.\ behavioral collapse} (Phase I,
discovery; day-1 population: 14 harnesses $\times$ 60 tasks, 91 pairs
including the \bare{} and \react{} baselines). Each point is a harness
pair; $x$ = token-sequence code similarity, $y$ = fraction of tasks with
differing correctness. Code similarity does not predict outcome
disagreement ($\rho=-0.01$); 36/91 pairs are identical on all 60 tasks.
Raw pairs: \texttt{fig1\_phenomenon\_data.json} in the code release.}
\label{fig:phenomenon}
\end{figure}

\subsection{Trace taxonomy: why code diversity fails to become behavioral
diversity}\label{sec:taxonomy}
Auditing execution traces (the sequence of LLM calls and artifact executions
per task) decomposes ``different code'' into three failure classes:
\begin{itemize}
\tightlist
\item \textbf{T1 --- mechanism described, not executed (6/12 candidates).}
The harness issues exactly one LLM call and never executes an artifact ---
the identical execution path as \bare{}. Docstrings claim retrieval,
self-checking, or token caps; none exists in the control flow.
\item \textbf{T2 --- feedback without a causal path (\react, 10/10 audited
tasks).} SQL is executed and errors fed back, yet the final SQL is
byte-identical to \bare's on every audited task.
\item \textbf{T3 --- broken multi-turn path (1/12).} The one candidate
implementing genuine multi-turn repair returns prose as final SQL,
converting bare-correct items into harness-wrong.
\end{itemize}
The taxonomy is actionable: T1 demands an execution-path acceptance gate,
T2 demands verified causal paths for feedback, T3 demands
executable-output validation --- all generation-protocol requirements.

\subsection{Human positive control: the headroom exists}\label{sec:positive}
To exclude the alternative that the domain simply lacks
instance-conditional signal, we wrote four harnesses by hand with forcibly
distinct execution paths (\texttt{hpc\_repair}: execute-and-repair;
\texttt{hpc\_vote3}: 3-vote on execution results; \texttt{hpc\_schema}:
schema-linking; \texttt{hpc\_hint}: hint-first reordering). Same frozen
target, tasks, and judge as the collapsed population: oracle headroom
\textbf{6.67\,pp} on 60 tasks and \textbf{7.28\,pp} {[}3.3, 10.0{]} on the
151-item set, with repair/schema/hint each producing repairs the others
miss. One disclosed caveat: the control was designed \emph{after}
inspecting day-1 traces, so its headroom is an existence proof, not a blind
estimate. The bare baseline itself flipped 3 of 60 items between the two
day-1 runs (API nondeterminism at temperature 0).

\subsection{Phase-I protocol probes and their limits}\label{sec:phase1probes}
Phase I also probed generation-side fixes --- a strategy-forced free-form protocol
(assigned mechanism, import/smoke acceptance) and a compiled Harness IR. Under the
strategy-forced protocol two of three builders produced populations whose point
estimates exceeded the pre-specified thresholds (Qwen 5.96, DeepSeek 7.95\,pp; the
old-protocol arm missed at 4.64\,pp), and per-strategy acceptance ordered with
builder identity (GLM 3/8, Qwen 7/8, DeepSeek 8/8, under the logged tool). These
observations motivate Phase II but are not confirmatory, for three disclosed reasons:
the design was two nested contrasts that cannot separate protocol from builder main
effects; the early runs lack per-attempt logs, so budget equality cannot be audited;
and acceptance was an import/smoke check, not the counterfactual conformance gate of
\S\ref{sec:phase2}. Full condition details are in Appendix~\ref{app:phase1-provenance}.

\paragraph{Generation-null failures (M0).} Re-running the original optimize loop
with Qwen and DeepSeek builders produced 40 candidates byte-identical to the
\bare{} baseline source. Identical code trivially yields identical outcomes, so this
run provides no evidence about behavioral collapse; it is reported as a distinct,
more primitive failure mode --- the pipeline produced no distinct population at all.
The collapse claims above rest exclusively on the genuinely distinct-source day-1
population.
} When we later re-ran the
original optimize loop with Qwen and DeepSeek builders, its 40 emitted
candidates proved byte-identical to the \bare{} baseline source. Identical
code trivially yields identical outcomes, so this run provides \emph{no}
evidence about behavioral collapse or builder effects; we report it as a
distinct, more primitive failure mode --- the pipeline produced no distinct
population at all. The collapse claims above rest exclusively on the
genuinely distinct-source day-1 population.
